\section{Why I Was Born}

\paragraph{Choosing One's Birth}

One often hears the lament, even from spiritual or religious people, ``Why should I be judged, when I didn't even choose to be born?'' \textbf{Valentin Tomberg}, in the \emph{Letter on Death}, formulates a response to that question:

\begin{quotex}
Birth can be either ``holy'' or ``natural'', i.e. it can either be an act of obedience to divine will or rather it can be effected as a consequence of a ``call from the earth''. A soul can be sent to the earth or it can be attracted by the earth.
\end{quotex}


\paragraph{Natural Birth}


\begin{quotex}
\textbf{Natural birth} is an event that is half-voluntary and half-involuntary, where the soul falls --- often without realising it --- into the sphere of terrestrial attraction, which bears it to birth, and thereby it is made little by little to forget its experiences above. Birth is then a forgetting of heaven and simultaneously a recall to the earth.
\end{quotex}

In this case, the person simply forgets his spiritual origin. Therefore, he becomes ignorant of his purpose and destiny on earth.

\paragraph{Holy Birth}


\begin{quotex}
\textbf{Holy birth} is the remembrance of the divine which is the force which accomplishes incarnation. It is not thanks to forgetting of the divine that the soul is then incarnated, but rather thanks to its remembering. The soul is incarnated in the state of habitual union with God. Then its will does not lose memory of the divine. This memory acts in it, imprinted as it is in the soul's will, during the whole terrestrial life which follows a holy birth.
\end{quotex}


In this case, the soul is united to the Will of God.

\paragraph{Mission and Tradition}


Following a holy birth, the soul is united with God in life, as well as in birth. Tomberg elucidates:

\begin{quotex}
The \textbf{true mission} is not what the human being proposes to do on the earth according to his tastes, his interests and even his ideals, but rather what God wants him to do. A true mission on the earth serves the cause of the ennoblement and spiritualisation of that which is, i.e. of what lives as tradition. It brings an impulse effecting the rejuvenation and intensification of tradition.
\end{quotex}


Arbitrary missions, on the contrary, aim at revolutionising the course of mankind's history and substituting specific innovations for what lives as tradition. Arbitrary missions have only contributed confusion to human history.

\paragraph{Fulfilling Destiny}


It is necessary to begin where you are, as a person, and then find your purpose in space and time:

\begin{itemize}


\item Start with your personal life, which is associated with a guardian angel.

\item To learn where your place in space is, you will develop a relationship with the archangel, who rules over your nation.

\item To put yourself in time, develop a relationship with the archai (or principality), the spirit of time that rules your era.
\end{itemize}


\paragraph{Persona}


People identify with their Persona, not the Self. The Persona is the appearance one presents to the world; the Persona can change, depending on situations, circumstances, or relationships. To the extent you identify with your Persona, you are presenting an inauthentic self.

The point of the personal stories is to demonstrate how malleable the Persona is. Most people think that their personality is unchangeable, but the true Self can mould it as needed. If you learn to do it appropriately and in the right circumstances, you may experience better life outcomes.

\paragraph{Platonism}


Before we begin our study of ur-platonism, let's review the philosophical argument:

\begin{itemize}


\item \textbf{Being}: We begin with a sense image of something


\item \textbf{Thinking}: We abstract from it to reach the form or idea of the thing

\item \textbf{Intuition}: Ultimately, the One is beyond being --- and beyond thought --- so it can only be known through intuition
\end{itemize}


We will see how Plato, Aristotle, Plotinus, and Thomas Aquinas understood these issues.

\paragraph{Spiritual Symbols}


But if we take that same schema and apply it to spiritual or hermetic symbols, we can gain great insights.

\begin{itemize}


\item \textbf{Imagination}: We start with a sacred symbol, visualized in the imagination

\item \textbf{Inspiration}: By meditating on the symbol, an understanding of the meaning of the symbol is gained

\item \textbf{Intuition}: Ultimately, one leaves both the symbol and the thought about it behind, and directly grasps the symbol through an intuitive gnosis
\end{itemize}


There is the temptation to assume that step (2) reveals the ``true meaning'' of the symbol. But there is really no difference between the symbol and its meaning. A temptation is to intellectualize the symbol, which ultimately does not explain it, but rather explains it away.

Tomberg offers this insight:

\begin{quotex}
Symbols are ``magic, mental, psychic and moral operations'' awakening new notions, ideas, sentiments and aspirations, which means to say that they require an activity more profound than that of study and intellectual explanation. It is therefore in a state of deep contemplation --- and always ever deeper --- that they should be approached.
\end{quotex}

These are the meeting notes for our esoteric classes on 23 November 2020. There are spots available for 2021. If you would like to be considered, drop me a note or leave a comment with your correct email address.


\flrightit{Posted on 2020-11-26 by Cologero}

